% =====================================================================================
\subsection{Zones from the inputs of a task: a bound that does not grow with later tasks}
\label{sec:bound}

Theorem~\ref{thm:writeonce} speaks about an input that is known to lie inside a zone. It does not
say how many inputs of a task do. For an unedited instruction the answer is all of them
(Proposition~\ref{prop:own}); for reworded instructions we could only measure it
(Section~\ref{sec:paraphrase}). There is a setting in which the share is guaranteed in advance, and
it needs no instruction at all: the zone is built from the task's own inputs, and its size is set
on inputs of the task that were held out. This subsection states the guarantee. In our tests it held, but with the features we
tried the zones of early tasks held the inputs of later tasks, which were then blocked
(Section~\ref{sec:noinstruction}); the recogniser of Section~\ref{sec:audit} avoids this. We keep
the statement because it is the one bound in this paper that needs no instruction and does not
grow with the number of later tasks.

\paragraph{Zones from the inputs of a task.} Let $u$ be a frozen feature of a whole input, as in
Section~\ref{sec:audit}; no part of the input is treated as an instruction. When task $k$ arrives with training inputs $D_k$, split
$D_k$ at random into fit inputs $D_k^{\mathrm{f}}$ and $m$ calibration inputs $c_1, \dots, c_m$.
Fix a score $a_k(x) \ge 0$ that says how unlike the fit inputs an input $x$ is. The score may depend
on $u$, on $D_k^{\mathrm{f}}$ and on anything frozen, but not on the calibration inputs. For a level
$\beta \in (0,1)$ (the level called $\alpha$ in Appendix~\ref{app:widen}; we write $\beta$ because $\alpha$ is the scale of the adapters in Eq.~\eqref{eq:layer}) let $q_k$ be the $\lceil (1-\beta)(m+1) \rceil$-th smallest of
$a_k(c_1), \dots, a_k(c_m)$, with $q_k = \infty$ if that index exceeds $m$, and let
\begin{equation}
  Z_k \;=\; \{\, x : a_k(x) \le q_k \,\} .
  \label{eq:inputzone}
\end{equation}
What defines the zone (the fit features and $q_k$) is stored and never changed. The switch is used
as before: an input goes to the oldest zone that holds it (Eq.~\eqref{eq:route}; an input outside every zone goes to the task with the smallest $a_k(x)/q_k$), and it is answered
with the adapters that existed when that task was learned. The adapter of task $k$ may be trained
on all of $D_k$; the zone does not depend on it.

\begin{theorem}[Bounded forgetting rate]
\label{thm:bound}
Assume the conditions of Theorem~\ref{thm:writeonce} (they apply to these zones by Remark~\ref{rem:otherzones}), with the zones of Eq.~\eqref{eq:inputzone},
and let $X$ be an input of task $k$ such that, given the fit inputs, $c_1, \dots, c_m, X$ are
exchangeable (for example, the training inputs of the task and $X$ are independent draws from one
distribution). Then for every choice of the later tasks, and of the way their adapters are trained,
\[
  \Pr\bigl[\ \text{there are } T' > T \ge k \text{ with } y_{T'}(X) \ne y_T(X) \ \bigr]
  \;\le\; \Pr\bigl[\, X \notin Z_k \,\bigr] \;\le\; \beta .
\]
\end{theorem}

\begin{proof}
Given the fit inputs, $a_k$ is a fixed function, so the $m+1$ numbers
$a_k(c_1), \dots, a_k(c_m), a_k(X)$ are exchangeable. The proof of
Proposition~\ref{prop:conformal} uses nothing else, and gives
$\Pr[a_k(X) \le q_k] \ge 1 - \beta$, that is, $\Pr[X \notin Z_k] \le \beta$. Now let $X \in Z_k$.
Then $X \in Z_1 \cup \dots \cup Z_k$, so $X$ is certified at time $k$ ($c_k(X) = 1$ in the notation of Eq.~\eqref{eq:certificate}), and Theorem~\ref{thm:writeonce}, applied at time $k$, gives
$y_{T}(X) = y_k(X)$ for every $T \ge k$ and every choice of the later tasks. So the answer to $X$
can differ between two times $T' > T \ge k$ only if $X \notin Z_k$.
\end{proof}

\begin{corollary}[A bound on the usual measure of forgetting]
\label{cor:bound}
Let $\mu$ be a fixed, deterministic score of an answer against its reference answers, with values
in $[0, M]$ (for Rouge-L in per cent, $M = 100$). Under the conditions of
Theorem~\ref{thm:bound}, for all $T' > T \ge k$ and all later tasks,
\[
  \bigl|\ \mathbb{E}\,\mu\bigl(y_T(X)\bigr) - \mathbb{E}\,\mu\bigl(y_{T'}(X)\bigr) \ \bigr|
  \;\le\; M \beta .
\]
In particular the expected forgetting of every task is at most $M\beta$ points, whatever is learned
later.
\end{corollary}

\begin{proof}
The two scores are equal unless the answer to $X$ changed between $T$ and $T'$, and then they differ
by at most $M$. By Theorem~\ref{thm:bound} this event has probability at most $\beta$.
\end{proof}

The bound covers an input drawn from the task's whole distribution, not only a marked subset of inputs; it uses no property of
the model, of the feature or of the later tasks; and it is set by a number that is chosen in
advance. Its price is paid when a task arrives, and it is visible there.

\begin{remark}[What the bound does not say]
\label{rem:boundlimits}
\emph{(i) It bounds change, not quality.} An input of task $k$ that lies in the zone of an older
task is answered by the adapters of that older task from the first moment on. It does not forget,
but task $k$ was never learned for it. The share of such inputs,
$b_k = \Pr[X \in Z_1 \cup \dots \cup Z_{k-1}]$, is not controlled by anything. It is measured on the
training inputs of the task when it arrives, and reported: the older zones do not depend on these
inputs, so for $n$ independent training inputs the observed share $\hat b_k$ satisfies
$|\hat b_k - b_k| \le \sqrt{\log(2/\delta)/(2n)}$ with probability at least $1 - \delta$
(Hoeffding's inequality). If two tasks have the same kind of input and differ only in what is asked,
$b_k$ is close to one: without an instruction nothing can tell them apart, and the report says so
before any training.
\emph{(ii) The probability is over the calibration inputs and the test input together.} For one
given set of calibration inputs the share of the task's inputs outside its zone can be larger than
$\beta$. For independent draws it is at most $\beta + \sqrt{\log(1/\delta)/(2m)}$ with probability
at least $1 - \delta$ over the calibration inputs (for $\delta \le 1/2$): by the one-sided
Dvoretzky--Kiefer--Wolfowitz inequality
the share of calibration scores at or below any threshold exceeds the true share by at most that
amount, and at least $\lceil (1-\beta)(m+1) \rceil$ of the $m$ calibration scores are at or below
$q_k$.
\emph{(iii) The test inputs have to be like the training inputs of the task.} An input from a
shifted distribution, for example a reworded one, has no guarantee from this theorem.
\emph{(iv) Something has to be stored for every task:} the fit features for a nearest-neighbour
score, or a mean and a radius for a distance to the mean.
\end{remark}

For the zones of Section~\ref{sec:zones}, built from the instruction, the same statement holds with
$\beta = 0$ for unedited inputs (Proposition~\ref{prop:own}). Theorem~\ref{thm:bound} is its
counterpart for inputs that share no fixed opening.
